\documentclass{article}
\usepackage{tcee-repaso}
\tcetema[Nueva teoría del comercio internacional]{3.B.6}{Teoría del comercio internacional (II). Nueva teoría del comercio internacional. Especial referencia a la competencia imperfecta, los rendimientos crecientes y la heterogeneidad empresarial.}
\tcefecha{05/11/2026}
\begin{document}
\cabecerarepaso
\begin{columnas}

% ==========================================================================
\bloque{Esquema del cante}
\esquemacante{%
\apartadocante{Introducción: «no vendas miel al colmenero»}{2,5}
\apartadocante{I. Competencia monopolística}{14}
\subcante{1. Antecedentes: sin ventaja comparativa estática}{1,5}
\subcante{2. Modelo de referencia (\autor{Krugman}, 1979)}{6}
\subcante{3. Costes iceberg y mercado doméstico (1980)}{2}
\subcante{4. Heterogeneidad: \autor{Melitz} (2003)}{4,5}
\apartadocante{II. Competencia oligopolista}{5}
\subcante{1. Cantidades: dumping recíproco}{3,5}
\subcante{2. Precios y diferenciación}{1,5}
\apartadocante{III. Otros desarrollos}{7}
\subcante{1. Economías externas y dinámicas. NGE}{5,5}
\subcante{2. Modelos de gravedad}{1,5}
\apartadocante{Conclusión}{1,5}}
\emph{Pizarra}: en una esquina, la tabla de modelos vacía, que se rellena con cada modelo; debajo, $p=\frac{\sigma}{\sigma-1}cw$ y el margen en cada modelo. En cada modelo, las 4 preguntas: por qué se comercia, qué produce cada país, a qué RRI y quién gana. Imprescindibles: PP-ZZ (autarquía y apertura), umbrales de \autor{Melitz} (autarquía, apertura con 5 grupos y ganadoras), reacciones del dumping, \emph{tomahawk}. Pegatinas: «cabeza de ratón\ldots» (\autor{Krugman}), «no todo el monte es orégano» (\autor{Melitz}), «pagar con la misma moneda» (dumping), «a quien madruga\ldots» (externas), «cuantos más, mejor» (C-P).

% ==========================================================================
\bloque[2,5]{Introducción}
\begin{itemize}
\item \textbf{Enganche}: «No vendas miel al colmenero»; pero España vende a Alemania $\approx$2.100 M€ de medicamentos y le compra $\approx$2.500 M€ (2025; GL = 0,92). La teoría tradicional \vertema{3.B.5} no lo explica: ¿por qué se vende miel al colmenero?
\item \textbf{Relevancia}: comercio mundial de bienes y servicios 34,65 billones \$ (2025, +7\,\%; OMC); intraindustrial 27\,\% (5 dígitos) y 44\,\% (3 dígitos) en 2006 (\autor{Brülhart}); Nobel a \autor{Krugman} (2008): patrones del comercio y localización; \autor{Melitz} (2003) abre la investigación con microdatos; política: cambian ganancias, ganadores y argumentos para intervenir \vertema{3.B.8}.
\item \textbf{Contextualización}: tradicional = clásicos (tecnología) y neoclásicos/H-O (dotaciones), con competencia perfecta, rendimientos constantes y bienes homogéneos. \emph{3 hechos sin explicar}: (1) mucho comercio intraindustrial ($\approx$½ dentro de Europa occ. y Norteamérica); (2) apertura de los ricos ↑ (31 $\to$ 60\,\% del PIB, 1970-2024) aunque sus dotaciones se parecen; (3) los desarrollados comercian entre sí (36\,\% del comercio de bienes, 2024). \emph{Nueva teoría} (70-80; \autor{Krugman}, \autor{Helpman}, \autor{Lancaster}, \autor{Brander}): \concepto{rendimientos crecientes} + \concepto{diferenciación} \vertema{3.A.18} \vertema{3.A.19}. \concepto{Equilibrio integrado} (\autor{Helpman}-\autor{Krugman}, 1985): ventaja comparativa $\to$ interindustrial; escala y diferenciación $\to$ intraindustrial. «Novísima» (2003): \concepto{heterogeneidad empresarial}. Antecedentes: \autor{Kravis} (disponibilidad, 1956), \autor{Posner}, \autor{Vernon}, \autor{Linder}, \autor{Falvey}.
\item \textbf{Problemática}: ¿por qué comercian países parecidos bienes parecidos? ¿Qué ganancias nuevas (variedad, escala, competencia, selección)? ¿Dónde se localiza la producción con rendimientos crecientes?
\item \textbf{Estructura}: \hyperref[rep:I]{I}, competencia monopolística (\autor{Krugman}, extensiones, \autor{Melitz}); \hyperref[rep:II]{II}, oligopolio (comportamiento estratégico); \hyperref[rep:III]{III}, economías externas y dinámicas, NGE y gravedad.
\end{itemize}

\textbf{Mapa de los modelos} (estructura de mercado × producto):
\begin{itemize}
\item \emph{Competencia perfecta}: homogéneo, \autor{Ricardo}, H-O-S \vertema{3.B.5}; vertical, neo-H-O de \autor{Falvey} (1981).
\item \emph{Competencia monopolística}: horizontal, \autor{Krugman} (1979, 1980), \autor{Lancaster} (1980), + heterogeneidad \autor{Melitz} (2003).
\item \emph{Oligopolio}: homogéneo, dumping recíproco (\autor{Brander}, 1981; \autor{Brander}-\autor{Krugman}, 1983); vertical, \autor{Shaked}-\autor{Sutton} (1984); horizontal, \autor{Eaton}-\autor{Kierzkowski} (1984).
\end{itemize}

% ==========================================================================
\bloque[12,5]{I. Competencia monopolística\label{rep:I}}
\begin{ideaclave}
Con preferencia por la variedad y economías de escala internas, 2 países idénticos comercian variedades distintas del mismo bien y ganan por la variedad; con elasticidad decreciente, también por competencia y escala; con costes de transporte, el mercado grande atrae a las empresas; con empresas heterogéneas, solo exportan las más productivas y la apertura reasigna hacia ellas.
\end{ideaclave}

\sub{I.2. Modelo de referencia (\autor{Krugman}, 1979)\label{rep:I.1}}
\begin{itemize}
\item \emph{Idea}: \concepto{rendimientos crecientes internos} (cada variedad en un solo sitio) + \concepto{demanda de bienes diferenciados} $\Rightarrow$ comercio entre países idénticos, intraindustrial por construcción. Precedentes sin modelo: \autor{Graham} (1923), \autor{Ohlin} (1933).
\item \emph{Supuestos}: lista de la competencia monopolística \vertema{3.A.18}: se levanta el producto homogéneo; muchas empresas, cada una monopolista de su variedad, toman $P$ dado; libre entrada. Preferencias \autor{Dixit}-\autor{Stiglitz} (1977):
\begin{formula}
\centering $U=\Big[\sum_{i=1}^{n}x_i^{\rho}\Big]^{1/\rho},\quad \rho\equiv\dfrac{\sigma-1}{\sigma},\ \sigma>1$
\end{formula}
\item[] \concepto{Preferencia por la variedad}: con $x_i=X/n$, $U=n^{1/(\sigma-1)}X$; más valiosa cuanto \emph{menor} $\sigma$. $\sigma>1$: demanda elástica y precio finito. Tecnología: $C(q)=w(F+cq)$, CMe ↓ $\Rightarrow$ \concepto{economías de escala internas} \vertema{3.A.12}. Sin ventaja comparativa; solo difieren $L_A$, $L_B$.
\item \emph{Demanda CES}: $x_i=\frac{w}{P}\big(\frac{p_i}{P}\big)^{-\sigma}$, $P\equiv[\sum p_j^{1-\sigma}]^{1/(1-\sigma)}$: ↑ con $w$ y $P$, ↓ con $p_i$; elasticidad $\sigma$.
\item \emph{Precio} (IMg = CMg), la regla del margen \vertema{3.A.17}:
\funcionclave{p^*=\frac{\sigma}{\sigma-1}\,c\,w\ \Longleftrightarrow\ \frac{p^*-cw}{p^*}=\frac{1}{\sigma}}
\item \emph{Libre entrada} ($p^*=$ CMe) y \emph{pleno empleo}:
\begin{formula}
\centering $q^*=\dfrac{(\sigma-1)F}{c},\qquad n^*=\dfrac{L}{\sigma F}$
\end{formula}
\item[] $q^*$ no depende de $L$; $n^*\propto L$.
\item \emph{Apertura} (sin costes) = sumar poblaciones: $p^*/w$ y $q^*$ iguales; cada país, sus $n_k^*$ variedades distintas («cabeza de ratón»); el consumidor tiene $n_A+n_B$. Importaciones:
\begin{formula}
\centering $M_A=\dfrac{w\,L_A L_B}{L_A+L_B}=M_B$
\end{formula}
\item[] equilibrado, intraindustrial, mayor cuanto más parecidos (gravedad sin distancia). \emph{Bienestar}: $w/P=\frac{\sigma-1}{\sigma c}\,n^{1/(\sigma-1)}$; 2 países iguales: ×$2^{1/(\sigma-1)}$: toda la ganancia es \emph{variedad}. Gráficamente: tangencia de \autor{Chamberlin}.
\item \emph{Elasticidad variable} (original de 1979): $U=\sum v(x_i)$, $|\varepsilon|$ ↓ con $x$. \concepto{Curva PP} (máx. beneficio), creciente: $p/w=c\,|\varepsilon|/(|\varepsilon|-1)$; \concepto{curva ZZ} (beneficio nulo), decreciente: $p/w=c+F/(Lx)$.
\item[] Apertura $L\to2L$: PP igual, ZZ baja $\Rightarrow$ (1) $x$ ↓; (2) $p/w$ ↓: \concepto{efecto procompetitivo}; (3) $q$ ↑: \concepto{efecto escala}; (4) $n$ por país ↓: \concepto{racionalización} (pero $2n'>n$; no dice cuáles cierran).
\end{itemize}
\graficorepaso[0.64\linewidth]{krugman-ppzz}{\autor{Krugman} (1979): la apertura baja ZZ; caen $x$ y el margen, sube $q$}
\begin{itemize}
\item \emph{Implicaciones}: se comercia por la estructura de mercado; qué variedades produce cada país, indeterminado; ganan los consumidores; sin conflicto entre factores (uno solo) $\Rightarrow$ costes de ajuste intraindustriales $<$ interindustriales; RRI = 1.
\item \emph{Evidencia}: ¿qué es una industria? CUCI de 1 a 5 dígitos, SA 6, NC 8; \concepto{sesgo de agregación} (\autor{Finger}, 1975): España-Francia 0,63 (CUCI 3), 0,48 (CUCI 5), 0,46 (SA 6); bilateral, no multilateral; saldo (\autor{Aquino}, 1978); horizontal/vertical por valores unitarios ($\pm$15\,\%); marginal (\autor{Brülhart}, 1994). \concepto{Índice de \autor{Grubel}-\autor{Lloyd}} (1971), 0 interindustrial, 1 intraindustrial:
\funcionclave{GL_k=1-\frac{|X_k-M_k|}{X_k+M_k}=\frac{2\min\{X_k,M_k\}}{X_k+M_k}}
\item[] País: media ponderada. \autor{Balassa} (1966): el complementario, media simple. Mayor entre desarrollados, grandes, de gustos parecidos y próximos (\autor{Greenaway}-\autor{Milner}, 1986); España $\approx$0,33 (5 díg.), Alemania 0,52. \autor{Broda}-\autor{Weinstein} (2006): variedades importadas EE.\,UU. ×3 (1972-2001), ganancia 2,6\,\% del PIB.
\item \emph{Valoración}: explica los 3 hechos y 3 ganancias (variedad, precios, escala); no sustituye a la tradicional (equilibrio integrado). Límites: patrón indeterminado; CES sin efecto procompetitivo; sin costes de comercio; empresas iguales.
\end{itemize}

\sub{I.3. Costes iceberg y efecto del mercado doméstico (\autor{Krugman}, 1980)\label{rep:I.2}}
\begin{itemize}
\item \concepto{Costes iceberg} (\autor{Samuelson}, 1952): enviar $\tau>1$ para que llegue 1; $p_x=\tau p^*$.
\item \concepto{Efecto del mercado doméstico}: el país con mayor mercado interior de un bien concentra una parte \emph{más que proporcional} de sus empresas y es \emph{exportador neto} (intuición de \autor{Linder}); con un solo sector de rendimientos crecientes, el grande paga \emph{salarios más altos}. Hacen falta escala \emph{y} costes de transporte: germen de la NGE.
\item \autor{Helpman}-\autor{Krugman} (1985), 2 factores: neto por dotaciones + variedades por escala. Multiproducto (\autor{Baldwin}-\autor{Ottaviano}, 2001): IDE recíproca sin eliminar el comercio intraindustrial \vertema{3.B.18}. Primer microfundamento de la gravedad.
\end{itemize}

\sub{I.4. Heterogeneidad empresarial: \autor{Melitz} (2003)\label{rep:I.3}}
\begin{itemize}
\item \emph{3 hechos} (\autor{Bernard}-\autor{Jensen}-\autor{Redding}-\autor{Schott}, 2007): (1) exportar es raro: 4\,\% de 5,5 M empresas EE.\,UU. (18\,\% en manufacturas); el 10\,\% mayor hace el 96\,\%; (2) las exportadoras son distintas: +97\,\% empleo, +108\,\% ventas, +11\,\% VAB/trabajador, +6\,\% salario; la productividad \emph{precede} a la exportación (\concepto{autoselección}); (3) abrirse reasigna (\autor{Trefler}, 2004).
\item \emph{Supuestos}: \autor{Krugman} CES (margen $\frac{\sigma}{\sigma-1}$ igual para todas) + productividad $\varphi$ (CMg $w/\varphi$) + coste hundido de entrada $wF^e$, luego se descubre $\varphi\sim G(\varphi)$. Costes: $C^d=w(F^d+q/\varphi)$; $C^x=w(F^x+\tau q_x/\varphi)$, con \concepto{coste fijo de exportar} $F^x$.
\item \emph{Autarquía}: $p(\varphi)=\frac{\sigma}{\sigma-1}\frac{w}{\varphi}$; beneficio operativo = ingreso$/\sigma$:
\begin{formula}
\centering $B^d(\varphi)=\dfrac{wL}{\sigma}\Big(\dfrac{\rho\varphi P}{w}\Big)^{\sigma-1}-wF^d$
\end{formula}
\item[] recta en $\varphi^{\sigma-1}$. Equilibrio: (1) beneficio nulo de la marginal, \concepto{productividad de corte} $B^d(\varphi^*_{\mathrm{aut}})=0$; (2) libre entrada: $[1-G(\varphi^*)]\bar B=wF^e$.
\item \emph{Comercio}: $B^x=\tau^{1-\sigma}\frac{wL}{\sigma}(\cdot)^{\sigma-1}-wF^x$, más abajo y menos pendiente. Umbrales:
\funcionclave{\varphi^*_x=\tau\Big(\frac{F^x}{F^d}\Big)^{\frac{1}{\sigma-1}}\varphi^*}
\item[] Si $\tau^{\sigma-1}F^x>F^d$, $\varphi^*_x>\varphi^*$. Al abrirse, las exportadoras crecen, ↑ demanda de trabajo, ↑ salario real ($P/w$ ↓): $B^d$ gira hacia abajo y $\varphi^*_{\mathrm{aut}}\to\varphi^*$ ↑. 5 grupos: (1) no producen; (2) \emph{cierran} (\concepto{efecto selección}); (3) solo en casa, venden y ganan menos; (4a) exportan pero ganaban más en autarquía ($\varphi^*_x\leq\varphi<\hat\varphi$); (4b) exportan y ganan ($\varphi\geq\hat\varphi$, donde $B^d+B^x=B^d_{\mathrm{aut}}$). «No todo el monte es orégano».
\end{itemize}
\graficorepaso[0.72\linewidth]{melitz-umbrales}{\autor{Melitz}: la apertura sube el umbral de supervivencia; solo exportan las de $\varphi\geq\varphi^*_x$}
\begin{itemize}
\item[] Canal: no los márgenes (constantes), sino la competencia por el trabajo.
\item \emph{Implicaciones}: \concepto{selección} y \concepto{reasignación}: la productividad media ↑ sin que ninguna empresa cambie la suya; nueva ganancia; salario real ↑, aunque $n$ disponible puede ↑ o ↓; costes de ajuste. RRI = 1. Política: reducir $F^x$ (promoción exterior) amplía la base exportadora.
\item \emph{Extensiones}: exportar o invertir (\autor{Helpman}-\autor{Melitz}-\autor{Yeaple}, 2004): cierran $<$ solo casa $<$ exportan $<$ filiales \vertema{3.B.18}; con ventaja comparativa (\autor{Bernard}-\autor{Redding}-\autor{Schott}, 2007): más selección en las de ventaja \vertema{3.B.5}; márgenes variables (\autor{Melitz}-\autor{Ottaviano}, 2008): vuelve el efecto procompetitivo; gravedad (\autor{Chaney}, 2008): \autor{Pareto}, la distancia actúa por el \emph{margen extensivo}; \autor{Eaton}-\autor{Kortum} (2002): ricardiano en competencia perfecta, con gravedad.
\item \emph{Evidencia}: ALC Canadá-EE.\,UU. (1989): escala débil, $\pm$10\,\% neto pequeño (\autor{Head}-\autor{Ries}, 1999); \emph{selección} fuerte: empleo $-12$\,\%, productividad de la industria $+15$\,\%, de las plantas $+14$\,\% (\autor{Trefler}, 2004). España 2025: 127.084 exportadoras, 46.229 regulares; las 1.000 primeras, 65,2\,\% \vertema{4.A.24}.
\item \emph{Valoración}: la empresa al centro; marco de la investigación de este siglo. Límites: productividad exógena; margen constante; ganancia total igual que en \autor{Krugman} (1980) y \autor{Eaton}-\autor{Kortum} (ACR, 2012): cambia de dónde sale y quién gana.
\end{itemize}

% ==========================================================================
\bloque[5,5]{II. Competencia oligopolista\label{rep:II}}
\begin{ideaclave}
La rivalidad estratégica crea comercio por sí sola: con cantidades y mercados segmentados se comercia el mismo bien en las 2 direcciones (dumping recíproco), con un efecto procompetitivo que puede no compensar el transporte. Con precios y producto homogéneo desaparece; vuelve con la diferenciación.
\end{ideaclave}

\sub{II.1. Cantidades: el dumping recíproco (\autor{Brander}, 1981; \autor{Brander}-\autor{Krugman}, 1983)\label{rep:II.1}}
\begin{itemize}
\item \emph{Idea}: ¿por qué se exporta e importa el \emph{mismo} bien homogéneo (cemento, acero)? Cada monopolista ve en el otro mercado una ocasión de vender con margen sin rebajar el precio en casa. «Pagar con la misma moneda».
\item \emph{Supuestos}: 1 empresa por país, bien homogéneo, CMg $c$; \concepto{mercados segmentados} (sin arbitraje); \autor{Cournot} en cada mercado \vertema{3.A.19}; iceberg $\tau\geq1$ (CMg extranjera $\tau c$); $p^A=a-bQ^A$.
\item \emph{Funciones de reacción} en el mercado $A$:
\funcionclave{R_A=\frac{a-c}{2b}-\frac{q_B^A}{2},\quad R_B=\frac{a-\tau c}{2b}-\frac{q_A^A}{2}}
\item[] Equilibrio de \autor{Cournot}-\autor{Nash}:
\begin{formula}
\centering $q_A^A=\dfrac{a-2c+\tau c}{3b},\ q_B^A=\dfrac{a+c-2\tau c}{3b},\ p^A=\dfrac{a+c+\tau c}{3}$
\end{formula}
\item[] $\tau=1$: \autor{Cournot} simétrico, $q^C=(a-c)/3b$. $\tau$ ↑: $R_B$ hacia el origen; $q_A^A$ ↑ ($c/3b$), $q_B^A$ ↓ el doble ($-2c/3b$), $Q$ ↓, $p$ ↑. Hay comercio si $\tau c<\frac{a+c}{2}=p^M$.
\item \emph{Implicaciones}: \concepto{dumping recíproco}: cada una vende fuera más barato, neto de transporte ($p^A/\tau$), porque allí su cuota es menor y su demanda percibida más elástica. Comercio intraindustrial de un bien idéntico. \emph{Bienestar ambiguo}: procompetitivo (rompe 2 monopolios) vs. despilfarro del transporte cruzado; con $\tau$ pequeño gana el primero; cerca del prohibitivo puede empeorar. Ej.: $a=10$, $b=1$, $c=2$: $W^{\mathrm{aut}}=24$, $W(\tau)=38-12\tau+\frac{22}{9}\tau^2$: mejora si $\tau<21/11\approx1{,}91$; sin comercio si $\tau\geq3$. Con libre entrada, domina el procompetitivo.
\item \emph{Extensión}: arancel o subvención a la exportación desplazan las reacciones y trasladan beneficios: \concepto{política comercial estratégica} (\autor{Brander}-\autor{Spencer}, 1985) \vertema{3.B.8}.
\item \emph{Valoración}: explica el comercio cruzado de homogéneos y que el comercio reduce el poder de mercado. Límite: depende de cantidades y segmentación.
\end{itemize}
\graficorepaso[0.62\linewidth]{dumping-reciproco}{Dumping recíproco: reacciones en el mercado $A$ sin costes de transporte y con ellos}

\sub{II.2. Competencia en precios y diferenciación\label{rep:II.2}}
\begin{itemize}
\item \emph{\autor{Bertrand}, homogéneo}: la nacional fija $\tau c$ (o $p^M$ si es menor) y se queda el mercado: \emph{no hay comercio}; solo la amenaza de importación limita su poder. El dumping exige cantidades o capacidad y luego precios (\autor{Kreps}-\autor{Scheinkman}, 1983). Con \autor{Bertrand}, la política óptima pasa de subvención a \emph{impuesto} a la exportación (\autor{Eaton}-\autor{Grossman}, 1986).
\item \emph{Horizontal} (\autor{Eaton}-\autor{Kierzkowski}, 1984): eligen variedad y precio; vuelve el comercio intraindustrial.
\item \emph{Vertical} (\autor{Shaked}-\autor{Sutton}, 1983, 1984): calidad con coste fijo (I+D) $\Rightarrow$ calidades muy distintas y \concepto{oligopolio natural}: el nº de calidades depende de la distribución de la renta, no del tamaño. Al integrarse: salen empresas, precios ↓, mismas calidades; las supervivientes venden en los 2. Contraste con \autor{Falvey}: allí, competencia perfecta y dotaciones.
\end{itemize}

% ==========================================================================
\bloque[8]{III. Otros desarrollos\label{rep:III}}
\begin{ideaclave}
Con economías externas, la industria se queda donde empezó antes; con dinámicas, la ventaja comparativa se crea con la experiencia. Con factores móviles, los vínculos concentran la industria cuando comerciar es barato, con equilibrios múltiples e histéresis: la integración puede crear centro y periferia entre regiones iguales. La gravedad tiene en la nueva teoría su fundamento.
\end{ideaclave}

\sub{III.1. Economías externas y dinámicas. Nueva geografía económica\label{rep:III.1}}
\begin{itemize}
\item \concepto{Economías de escala externas} (\autor{Marshall}, 1890): CMe de cada empresa ↓ con la producción de \emph{la industria} en un lugar; compatibles con competencia perfecta. 3 fuentes: proveedores especializados, mercado de trabajo especializado, desbordamientos de conocimiento \vertema{3.A.12}.
\item \emph{Modelo} (los botones, KOM cap.~7): país 2 con CMe menor, pero la industria empezó en el 1. Al abrirse, se concentra en un país y el precio ↓ en los 2; el 2 no entra porque su $c_0$ supera el precio del 1. «A quien madruga, Dios le ayuda».
\end{itemize}
\graficorepaso[0.56\linewidth]{economias-externas}{Economías externas: la industria que empieza antes se queda el mercado aunque tenga costes mayores}
\begin{itemize}
\item[] Implicaciones: el patrón depende de la historia; el comercio puede empeorar a un país (\autor{Graham}, 1923; \autor{Ethier}, 1982) $\Rightarrow$ \emph{industria naciente} \vertema{3.B.8} (saber cuál, cuánto tiempo; captura). Externas internacionales (\autor{Ethier}, 1982): comercio de componentes sin concentrar; FPP convexa \vertema{3.A.11}.
\item \concepto{Economías de escala dinámicas}: coste ↓ con la producción \emph{acumulada} (curva de aprendizaje; \autor{Arrow}, 1962): ventaja comparativa \emph{creada}; un choque temporal (sobrevaluación) cambia la especialización (\autor{Krugman}, 1987) \vertema{3.B.9}.
\item \textbf{NGE}: \concepto{primera naturaleza} (recursos, ríos, puertos) vs. \concepto{segunda naturaleza} (decisiones que se refuerzan). \concepto{Nueva geografía económica} (\autor{Fujita}, 1988; \autor{Krugman}, 1991; \autor{Venables}, 1996): D-S + rendimientos crecientes + iceberg + factores móviles. Aglomeración sin diferencias exógenas.
\item \emph{Centro-periferia} (\autor{Krugman}, 1991): 2 regiones idénticas; agricultura homogénea, rendimientos constantes, sin transporte, numerario, trabajadores \emph{inmóviles} $(1-\mu)/2$ por región; manufacturas D-S, iceberg, trabajadores \emph{móviles} ($\lambda_i$). Preferencias:
\begin{formula}
\centering $U=X^{1-\mu}V^{\mu},\qquad \omega_i=w_iP_i^{-\mu}$
\end{formula}
\item[] Corto plazo: $n_i=\lambda_i\mu/(\sigma F)$. Dinámica, hacia el mayor salario real:
\funcionclave{\dot\lambda_A=\omega_A(\lambda_A)-\omega_B(\lambda_A)}
\item[] Equilibrios: \concepto{disperso} ($\lambda_A=\frac12$) o \concepto{centro-periferia} ($\lambda_A=0$ o $1$).
\item \emph{3 fuerzas} ante $\lambda_A$ ↑: \concepto{vínculo de demanda}: $Y_A$ ↑ $\Rightarrow$ $w_A$ ↑ (mercado doméstico), aglomera; \concepto{vínculo del coste de la vida}: $n_A$ ↑ $\Rightarrow$ $P_A$ ↓, aglomera; \concepto{competencia local} (\emph{market crowding}) + demanda agrícola inmóvil: dispersa. Los 2 primeros: \concepto{causación circular}. «Cuantos más, mejor».
\item \concepto{Libertad comercial} $\phi_\tau=\tau^{1-\sigma}\in(0,1]$: costes altos, solo disperso; intermedios, disperso y C-P estables (la historia decide); bajos, C-P.
\end{itemize}
\graficorepaso[0.68\linewidth]{tomahawk}{\emph{Tomahawk}: estables en continuo; disperso hasta $\phi_\tau^{\mathrm r}$, C-P desde $\phi_\tau^{\mathrm s}$; entre ambos, histéresis}
\begin{itemize}
\item[] \concepto{Punto de ruptura} $\phi^{\mathrm r}_\tau$ (\emph{break}), \concepto{de sostenibilidad} $\phi^{\mathrm s}_\tau<\phi^{\mathrm r}_\tau$ (\emph{sustain}); entre ellos, \concepto{histéresis}. \concepto{«No agujero negro»}: $\mu<\rho=(\sigma-1)/\sigma$ (si no, concentra con cualquier coste). La integración \emph{puede concentrar} \vertema{3.B.10} $\Rightarrow$ política regional \vertema{3.B.40}.
\item \emph{Extensiones}: vínculos verticales (\autor{Venables}, 1996): aglomera sin migración (mejor para países); geografía y crecimiento (\autor{Baldwin}-\autor{Forslid}, 2000) \vertema{3.A.44} \vertema{3.B.9}.
\item \emph{Evidencia}: división de Alemania: ciudades occidentales cerca de la frontera crecen menos (\autor{Redding}-\autor{Sturm}, 2008); ciudades japonesas bombardeadas recuperan su tamaño: pesan los fundamentos (\autor{Davis}-\autor{Weinstein}, 2002).
\item \emph{Valoración}: fuerzas endógenas, aglomeración, equilibrios múltiples, histéresis. Límites (\autor{Neary}, 2001): muy estilizado y numérico; depende de formas funcionales; poca migración entre países. \autor{Krugman} (1998): más éxito en las preguntas que en las respuestas.
\end{itemize}

\sub{III.2. Modelos de gravedad\label{rep:III.2}}
\begin{itemize}
\item \concepto{Ecuación de gravedad} (\autor{Tinbergen}, 1962; \autor{Isard}, 1954), la regularidad más robusta:
\begin{formula}
\centering $\ln X_{ij}=\beta_0+\beta_1\ln Y_i+\beta_2\ln Y_j-\beta_3\ln d_{ij}+u_{ij}$
\end{formula}
\item[] Elasticidades $\approx$1 (\autor{Head}-\autor{Mayer}, 2014). Nació sin teoría; la fundamenta la preferencia por la variedad; \concepto{resistencias multilaterales} (\autor{Anderson}-\autor{van Wincoop}, 2003); margen extensivo (\autor{Chaney}, 2008). Física: $F=GM_1M_2/d^2$; en forma multiplicativa, $X_{ij}=e^{\beta_0}Y_i^{\beta_1}Y_j^{\beta_2}/d_{ij}^{\beta_3}$. Fundamento: \autor{Anderson} (1979), \autor{Krugman}; estructural $X_{ij}=\frac{Y_iY_j}{Y^W}\big(\frac{\tau_{ij}}{\Pi_iP_j}\big)^{1-\sigma}$. \concepto{Efecto frontera} (\autor{McCallum}, 1995: $\approx$20 veces entre provincias canadienses): administrativo, geográfico, cultural, \emph{home bias}. Límite: demasiados fundamentos. Estimación y acuerdos \vertema{3.B.10}.
\end{itemize}

% ==========================================================================
\bloque[1,5]{Conclusión}
\begin{itemize}
\item \textbf{Recapitulación}: monopolística: escala + variedad bastan; precios ↓ si la elasticidad no es constante; con heterogeneidad, exportan las mejores. Oligopolio: se comercia el mismo bien. Externas y movilidad: historia y geografía. Respuestas: estructura de mercado; 4 ganancias (variedad, escala, competencia, reasignación); costes de comerciar. La tradicional no queda superada, sino completada.
\item \textbf{Relevancia}: ganancias nuevas, pero costes de ajuste en las que cierran: Canadá-EE.\,UU., empleo $-12$\,\% a CP, productividad $+14$-15\,\% a LP (\autor{Trefler}).
\item \textbf{Extensiones}: política estratégica \vertema{3.B.8}; integración \vertema{3.B.10}; comercio y crecimiento (\autor{Grossman}-\autor{Helpman}, 1991) \vertema{3.B.9}; cadenas globales de valor \vertema{3.B.18}.
\item \textbf{Opinión}: ACR (2012): la ganancia agregada sale de 2 datos (parte del gasto nacional, elasticidad del comercio); los modelos nuevos cambian \emph{cómo} y \emph{quién}, no \emph{cuánto} $\Rightarrow$ apertura + políticas para cruzar el umbral exportador y reasignar trabajadores.
\item \textbf{Idea final}: vender miel al colmenero tiene sentido si la miel es distinta, si producirla en grande es más barato y si quien la vende es el mejor colmenero: la diferencia entre países es causa y también \emph{resultado} del comercio.
\end{itemize}

\textbf{Tradicional vs. nueva}: países distintos vs. también iguales con escala; ventaja comparativa vs. además historia; trigo vs. aviones; protección: reduce el bienestar vs. puede ayudar con escala (si todos, fragmenta).

% ==========================================================================
\bloque{I.1. Antecedentes: comercio sin ventaja comparativa estática (nexo)}
\begin{itemize}
\item Nexo: los neotecnológicos y \autor{Linder}, comercio sin ventaja comparativa estática; \autor{Falvey}, intraindustrial en competencia perfecta. \autor{Kravis} (1956), \concepto{disponibilidad} (absoluta o relativa: innovación, diferenciación).
\item \autor{Posner} (1961), \concepto{desfase tecnológico}: de imitación vs. de demanda; si el de demanda es más corto, el innovador exporta: $t_1$ nace la demanda en $B$, crece hasta $t_2$ (imita), se anula en $t_3$.
\item \autor{Vernon} (1966), \concepto{ciclo del producto}: introducción (innovador), madurez (otros desarrollados, IDE), estandarización (salarios bajos exportan): ventaja comparativa dinámica.
\item \autor{Linder} (1961), \concepto{demanda representativa} y \concepto{demandas solapadas}: la renta pc fija la gama de calidades; más comercio entre rentas parecidas; evidencia sectorial a favor (\autor{Hallak}, 2010).
\item \autor{Falvey} (1981), neo-H-O: $p_i(\alpha)=w_i+\alpha r_i$; \concepto{calidad frontera} $\alpha_0=(w_A-w_B)/(r_B-r_A)$: el abundante en capital exporta las altas. Solo vertical.
\end{itemize}
\pareja{posner}{\autor{Posner}: exportaciones netas del innovador}{falvey}{\autor{Falvey}: la calidad frontera $\alpha_0$}

% ==========================================================================
\bloque{Autores y fechas}
\hito{1890}{\autor{Marshall}: economías externas}
\hito{1961}{\autor{Posner}: desfase; \autor{Linder}: demandas solapadas}
\hito{1962}{\autor{Tinbergen}: gravedad; \autor{Arrow}: aprendizaje}
\hito{1966}{\autor{Vernon}: ciclo del producto; \autor{Balassa}: CEE}
\hito{1971}{\autor{Grubel}-\autor{Lloyd}: índice}
\hito{1977}{\autor{Dixit}-\autor{Stiglitz}: competencia monopolística}
\hito{1979}{\autor{Krugman}: rendimientos crecientes y comercio}
\hito{1980}{\autor{Krugman}: mercado doméstico; \autor{Lancaster}}
\hito{1981}{\autor{Brander}: dumping; \autor{Falvey}: calidades}
\hito{1983}{\autor{Brander}-\autor{Krugman}: dumping recíproco}
\hito{1984}{\autor{Shaked}-\autor{Sutton}; \autor{Eaton}-\autor{Kierzkowski}}
\hito{1985}{\autor{Helpman}-\autor{Krugman}: equilibrio integrado; \autor{Brander}-\autor{Spencer}}
\hito{1991}{\autor{Krugman}: centro-periferia}
\hito{2002}{\autor{Eaton}-\autor{Kortum}; \autor{Davis}-\autor{Weinstein}}
\hito{2003}{\autor{Melitz}; \autor{Anderson}-\autor{van Wincoop}}
\hito{2004}{\autor{Trefler}; \autor{Helpman}-\autor{Melitz}-\autor{Yeaple}}
\hito{2008}{Nobel a \autor{Krugman}; \autor{Chaney}; \autor{Melitz}-\autor{Ottaviano}}
\hito{2012}{\autor{Arkolakis}-\autor{Costinot}-\autor{Rodríguez-Clare}}

% ==========================================================================
\bloque{Datos clave}
\begin{itemize}
\item Comercio mundial bienes y servicios: \dato{34,65 billones \$, +7\,\%}{OMC, 2025}.
\item Intraindustrial mundial: \dato{27\,\% (5 díg.), 44\,\% (3 díg.)}{Brülhart, 2006}.
\item Apertura renta alta: \dato{31 $\to$ 60\,\% del PIB}{BM, 1970-2024}; comercio entre desarrollados \dato{36\,\%}{UNCTAD, 2024}.
\item Medicamentos España-Alemania: \dato{2.100 vs. 2.500 M€, GL 0,92}{Eurostat, 2025}.
\item EE.\,UU.: exportan \dato{4\,\% de las empresas}{BJRS, 2000}; 10\,\% mayor \dato{96\,\% de la exportación}{BJRS}.
\item Variedad importada EE.\,UU.: \dato{×3, 2,6\,\% del PIB}{Broda-Weinstein, 1972-2001}.
\item España: \dato{127.084 exportadoras, 46.229 regulares}{MINECO, 2025}; 1.000 primeras \dato{65,2\,\%}{MINECO}; $\geq$250 empleados \dato{1,9\,\% de las empresas, 57\,\% del valor}{Eurostat, 2024}.
\end{itemize}

% ==========================================================================
\bloque{Preguntas del tribunal}
\begin{enumerate}
\item \textbf{¿Por qué comercian 2 países idénticos?} Preferencia por la variedad + escala interna: cada variedad en un solo sitio y se demandan todas (\autor{Krugman}, 1979).
\item \textbf{¿Qué añade la elasticidad variable a la CES?} Con CES solo variedad; con $|\varepsilon|$ decreciente, efecto procompetitivo, escala y racionalización (PP-ZZ).
\item \textbf{¿Qué es el efecto del mercado doméstico?} El país de mercado mayor concentra más que proporcionalmente las empresas y exporta en neto; exige escala y transporte.
\item \textbf{¿Por qué sube la productividad con la apertura en \autor{Melitz}?} Selección y reasignación hacia las más productivas; ninguna empresa cambia su $\varphi$; el canal es la competencia por el trabajo.
\item \textbf{¿Por qué $\varphi^*_x>\varphi^*$?} Por $F^x$ y $\tau$: si $\tau^{\sigma-1}F^x>F^d$.
\item \textbf{¿El dumping recíproco mejora el bienestar?} Ambiguo: procompetitivo vs. transporte cruzado; con libre entrada, sí.
\item \textbf{¿Y con \autor{Bertrand}?} Sin comercio de homogéneos; la política óptima pasa a impuesto a la exportación.
\item \textbf{¿Qué es el \emph{tomahawk}?} Equilibrios según $\phi_\tau$: disperso hasta el punto de ruptura, C-P desde el de sostenibilidad; histéresis entre ambos.
\item \textbf{¿Ganan más los modelos nuevos?} No necesariamente (ACR, 2012): cambian el origen y el reparto de la ganancia.
\item \textbf{¿Economías externas justifican proteger?} Industria naciente: la ventaja se queda donde se empezó; pero hay que saber cuál, cuánto tiempo, y riesgo de captura \vertema{3.B.8}.
\end{enumerate}

\begin{trampa}
\begin{itemize}
\item Variedad más valiosa con $\sigma$ \emph{baja}; margen $\frac{\sigma}{\sigma-1}$ ↓ con $\sigma$; $q^*=(\sigma-1)F/c$ ↑ con $\sigma$ y \emph{no depende} de $L$.
\item CES: apertura sin cambio de precios ni escala; PP-ZZ: baja ZZ, PP no se mueve.
\item Racionalización: menos empresas \emph{por país}, pero más variedades disponibles ($2n'>n$).
\item GL: 0 interindustrial, 1 intraindustrial; \autor{Balassa} (1966) mide lo contrario, con media simple.
\item \autor{Melitz}: margen \emph{constante}; la selección no viene de los márgenes sino del salario real; productividad exógena.
\item Dumping: con $\tau$ ↑, $q_B^A$ cae el \emph{doble} de lo que sube $q_A^A$.
\item Economías externas: compatibles con competencia perfecta; la industria no va al de menores costes.
\item $\phi_\tau=\tau^{1-\sigma}$: 1 sin costes; $\phi^{\mathrm s}_\tau<\phi^{\mathrm r}_\tau$. «No agujero negro»: $\mu<\rho$. $\mu$ = gasto en manufacturas, no margen.
\item \autor{Falvey}: competencia perfecta y H-O (vertical); \autor{Shaked}-\autor{Sutton}: oligopolio con costes fijos.
\item \autor{Eaton}-\autor{Kortum}: ricardiano en competencia \emph{perfecta}.
\item Efecto del mercado doméstico: \emph{más que proporcional}; exige escala \emph{y} costes de transporte; el país grande paga salarios \emph{más altos}.
\item \autor{Bertrand} con homogéneo: \emph{no} hay comercio; política estratégica óptima: \emph{impuesto}, no subvención.
\end{itemize}
\end{trampa}

\end{columnas}
\end{document}
